% 主控制状态机 FSM 图
% 使用方法: 在主文档中 % 主控制状态机 FSM 图
% 使用方法: 在主文档中 % 主控制状态机 FSM 图
% 使用方法: 在主文档中 % 主控制状态机 FSM 图
% 使用方法: 在主文档中 \input{figures/main_fsm.tex}
% 字体: 中文使用宋体（继承文档设置），英文使用 Times New Roman（newtxtext）

\begin{tikzpicture}[
    every node/.append style={font=\rmfamily},  % 确保使用衬线字体（Times/宋体）
    state/.style={draw, thick, rounded corners=5pt, minimum width=2.6cm, minimum height=0.9cm, align=center, font=\footnotesize},
    state_idle/.style={state, fill=gray!20},
    state_load/.style={state, fill=blue!15},
    state_phase2/.style={state, fill=green!15},
    state_phase3/.style={state, fill=orange!15},
    state_write/.style={state, fill=purple!15},
    state_done/.style={state, fill=yellow!20},
    arrow/.style={->, >=Stealth, thick},
    condition/.style={font=\tiny, align=center},
    loop_box/.style={draw, dashed, thick, rounded corners=8pt, fill=white, fill opacity=0.7},
]

% ========== 状态节点 ==========
\node[state_idle] (idle) at (0, 0) {IDLE};

\node[state_load] (load_param) at (0, -1.8) {LOAD\_PARAMS\\读取参数};

\node[state_load] (load_input) at (0, -3.6) {LOAD\_INPUTS\\读取SOB/观测/路标};

\node[state_phase2] (p2_init) at (0, -5.4) {PHASE2\_INIT\\初始化循环计数器};

% Phase 2 循环框
\node[loop_box, minimum width=6.5cm, minimum height=4.2cm] (p2_loop_box) at (4.5, -7.8) {};
\node[anchor=north, font=\footnotesize\bfseries] at (4.5, -5.8) {PHASE2\_LOOP};

\node[state_phase2] (p2_check) at (4.5, -6.8) {SOB 判定\\(state, cam)有效?};
\node[state_phase2] (coord_trans) at (2.5, -8.2) {COORD\\TRANSFORM};
\node[state_phase2] (residual) at (4.5, -8.2) {RESIDUAL\\CALC};
\node[state_phase2] (jacobian) at (6.5, -8.2) {JACOBIAN\\CALC};
\node[state_phase2] (hessian) at (4.5, -9.4) {HESSIAN\\ACCUM};

\node[state_phase3] (p3_init) at (0, -11) {PHASE3\_INIT\\初始化Schur循环};

% Phase 3 循环框
\node[loop_box, minimum width=6.5cm, minimum height=3.5cm] (p3_loop_box) at (4.5, -12.8) {};
\node[anchor=north, font=\footnotesize\bfseries] at (4.5, -10.9) {PHASE3\_LOOP};

\node[state_phase3] (v_inv) at (2.5, -12) {$\mathbf{H}_{mm}^{-1}$计算\\(j==0)};
\node[state_phase3] (wvinv) at (6.5, -12) {$\mathbf{H}_{pm}\mathbf{H}_{mm}^{-1}$\\计算};
\node[state_phase3] (schur_update) at (4.5, -13.5) {Schur 更新\\$\mathbf{H}_{pp}, \mathbf{b}_p$};

\node[state_write] (write_out) at (0, -15) {WRITE\_OUTPUT\\DMA写回结果};

\node[state_done] (complete) at (0, -16.8) {COMPLETE\\置位ICR.COMPLETE};

% ========== 状态转移 ==========
% IDLE -> LOAD_PARAMS
\draw[arrow] (idle) -- node[right, condition] {IP\_ENABLE\\$\land$ START\_H} (load_param);

% LOAD_PARAMS -> LOAD_INPUTS
\draw[arrow] (load_param) -- node[right, condition] {读取完成} (load_input);

% LOAD_INPUTS -> PHASE2_INIT
\draw[arrow] (load_input) -- node[right, condition] {加载完成} (p2_init);

% PHASE2_INIT -> PHASE2_LOOP
\draw[arrow] (p2_init) -- ++(2, 0) |- (p2_check);

% Phase 2 内部流程
\draw[arrow] (p2_check) -- node[left, condition] {有效} (coord_trans);
\draw[arrow] (coord_trans) -- (residual);
\draw[arrow] (residual) -- (jacobian);
\draw[arrow] (jacobian) -- (hessian);

% Phase 2 跳过路径
\draw[arrow, dashed] (p2_check.east) -- ++(1.5, 0) node[right, condition] {无效\\跳过} |- (hessian.east);

% Phase 2 循环返回
\draw[arrow] (hessian.south) -- ++(0, -0.3) -- ++(3.5, 0) -- ++(0, 3.0) -- ++(-3.5, 0) -- (p2_check.east);
\node[condition] at (8.5, -8) {cam++\\state++\\pt++};

% PHASE2 -> PHASE3_INIT
\draw[arrow] (hessian.west) -- ++(-3, 0) |- node[left, condition, pos=0.75] {pt $\geq$ num\_pts} (p3_init);

% PHASE3_INIT -> PHASE3_LOOP
\draw[arrow] (p3_init) -- ++(2, 0) |- (v_inv);

% Phase 3 内部流程
\draw[arrow] (v_inv) -- (wvinv);
\draw[arrow] (wvinv) -- (schur_update);

% Phase 3 循环返回
\draw[arrow] (schur_update.south) -- ++(0, -0.3) -- ++(3.5, 0) -- ++(0, 2.0) -- ++(-5.5, 0) -- (v_inv.west);
\node[condition] at (8.5, -12.5) {pt++\\j++};

% PHASE3 -> WRITE_OUTPUT
\draw[arrow] (schur_update.west) -- ++(-3, 0) |- node[left, condition, pos=0.8] {j $\geq$ state\_size} (write_out);

% WRITE_OUTPUT -> COMPLETE
\draw[arrow] (write_out) -- node[right, condition] {写入完成} (complete);

% COMPLETE -> IDLE (返回)
\draw[arrow] (complete.west) -- ++(-2, 0) -- ++(0, 16.8) -- (idle.west);

\end{tikzpicture}

% 字体: 中文使用宋体（继承文档设置），英文使用 Times New Roman（newtxtext）

\begin{tikzpicture}[
    every node/.append style={font=\rmfamily},  % 确保使用衬线字体（Times/宋体）
    state/.style={draw, thick, rounded corners=5pt, minimum width=2.6cm, minimum height=0.9cm, align=center, font=\footnotesize},
    state_idle/.style={state, fill=gray!20},
    state_load/.style={state, fill=blue!15},
    state_phase2/.style={state, fill=green!15},
    state_phase3/.style={state, fill=orange!15},
    state_write/.style={state, fill=purple!15},
    state_done/.style={state, fill=yellow!20},
    arrow/.style={->, >=Stealth, thick},
    condition/.style={font=\tiny, align=center},
    loop_box/.style={draw, dashed, thick, rounded corners=8pt, fill=white, fill opacity=0.7},
]

% ========== 状态节点 ==========
\node[state_idle] (idle) at (0, 0) {IDLE};

\node[state_load] (load_param) at (0, -1.8) {LOAD\_PARAMS\\读取参数};

\node[state_load] (load_input) at (0, -3.6) {LOAD\_INPUTS\\读取SOB/观测/路标};

\node[state_phase2] (p2_init) at (0, -5.4) {PHASE2\_INIT\\初始化循环计数器};

% Phase 2 循环框
\node[loop_box, minimum width=6.5cm, minimum height=4.2cm] (p2_loop_box) at (4.5, -7.8) {};
\node[anchor=north, font=\footnotesize\bfseries] at (4.5, -5.8) {PHASE2\_LOOP};

\node[state_phase2] (p2_check) at (4.5, -6.8) {SOB 判定\\(state, cam)有效?};
\node[state_phase2] (coord_trans) at (2.5, -8.2) {COORD\\TRANSFORM};
\node[state_phase2] (residual) at (4.5, -8.2) {RESIDUAL\\CALC};
\node[state_phase2] (jacobian) at (6.5, -8.2) {JACOBIAN\\CALC};
\node[state_phase2] (hessian) at (4.5, -9.4) {HESSIAN\\ACCUM};

\node[state_phase3] (p3_init) at (0, -11) {PHASE3\_INIT\\初始化Schur循环};

% Phase 3 循环框
\node[loop_box, minimum width=6.5cm, minimum height=3.5cm] (p3_loop_box) at (4.5, -12.8) {};
\node[anchor=north, font=\footnotesize\bfseries] at (4.5, -10.9) {PHASE3\_LOOP};

\node[state_phase3] (v_inv) at (2.5, -12) {$\mathbf{H}_{mm}^{-1}$计算\\(j==0)};
\node[state_phase3] (wvinv) at (6.5, -12) {$\mathbf{H}_{pm}\mathbf{H}_{mm}^{-1}$\\计算};
\node[state_phase3] (schur_update) at (4.5, -13.5) {Schur 更新\\$\mathbf{H}_{pp}, \mathbf{b}_p$};

\node[state_write] (write_out) at (0, -15) {WRITE\_OUTPUT\\DMA写回结果};

\node[state_done] (complete) at (0, -16.8) {COMPLETE\\置位ICR.COMPLETE};

% ========== 状态转移 ==========
% IDLE -> LOAD_PARAMS
\draw[arrow] (idle) -- node[right, condition] {IP\_ENABLE\\$\land$ START\_H} (load_param);

% LOAD_PARAMS -> LOAD_INPUTS
\draw[arrow] (load_param) -- node[right, condition] {读取完成} (load_input);

% LOAD_INPUTS -> PHASE2_INIT
\draw[arrow] (load_input) -- node[right, condition] {加载完成} (p2_init);

% PHASE2_INIT -> PHASE2_LOOP
\draw[arrow] (p2_init) -- ++(2, 0) |- (p2_check);

% Phase 2 内部流程
\draw[arrow] (p2_check) -- node[left, condition] {有效} (coord_trans);
\draw[arrow] (coord_trans) -- (residual);
\draw[arrow] (residual) -- (jacobian);
\draw[arrow] (jacobian) -- (hessian);

% Phase 2 跳过路径
\draw[arrow, dashed] (p2_check.east) -- ++(1.5, 0) node[right, condition] {无效\\跳过} |- (hessian.east);

% Phase 2 循环返回
\draw[arrow] (hessian.south) -- ++(0, -0.3) -- ++(3.5, 0) -- ++(0, 3.0) -- ++(-3.5, 0) -- (p2_check.east);
\node[condition] at (8.5, -8) {cam++\\state++\\pt++};

% PHASE2 -> PHASE3_INIT
\draw[arrow] (hessian.west) -- ++(-3, 0) |- node[left, condition, pos=0.75] {pt $\geq$ num\_pts} (p3_init);

% PHASE3_INIT -> PHASE3_LOOP
\draw[arrow] (p3_init) -- ++(2, 0) |- (v_inv);

% Phase 3 内部流程
\draw[arrow] (v_inv) -- (wvinv);
\draw[arrow] (wvinv) -- (schur_update);

% Phase 3 循环返回
\draw[arrow] (schur_update.south) -- ++(0, -0.3) -- ++(3.5, 0) -- ++(0, 2.0) -- ++(-5.5, 0) -- (v_inv.west);
\node[condition] at (8.5, -12.5) {pt++\\j++};

% PHASE3 -> WRITE_OUTPUT
\draw[arrow] (schur_update.west) -- ++(-3, 0) |- node[left, condition, pos=0.8] {j $\geq$ state\_size} (write_out);

% WRITE_OUTPUT -> COMPLETE
\draw[arrow] (write_out) -- node[right, condition] {写入完成} (complete);

% COMPLETE -> IDLE (返回)
\draw[arrow] (complete.west) -- ++(-2, 0) -- ++(0, 16.8) -- (idle.west);

\end{tikzpicture}

% 字体: 中文使用宋体（继承文档设置），英文使用 Times New Roman（newtxtext）

\begin{tikzpicture}[
    every node/.append style={font=\rmfamily},  % 确保使用衬线字体（Times/宋体）
    state/.style={draw, thick, rounded corners=5pt, minimum width=2.6cm, minimum height=0.9cm, align=center, font=\footnotesize},
    state_idle/.style={state, fill=gray!20},
    state_load/.style={state, fill=blue!15},
    state_phase2/.style={state, fill=green!15},
    state_phase3/.style={state, fill=orange!15},
    state_write/.style={state, fill=purple!15},
    state_done/.style={state, fill=yellow!20},
    arrow/.style={->, >=Stealth, thick},
    condition/.style={font=\tiny, align=center},
    loop_box/.style={draw, dashed, thick, rounded corners=8pt, fill=white, fill opacity=0.7},
]

% ========== 状态节点 ==========
\node[state_idle] (idle) at (0, 0) {IDLE};

\node[state_load] (load_param) at (0, -1.8) {LOAD\_PARAMS\\读取参数};

\node[state_load] (load_input) at (0, -3.6) {LOAD\_INPUTS\\读取SOB/观测/路标};

\node[state_phase2] (p2_init) at (0, -5.4) {PHASE2\_INIT\\初始化循环计数器};

% Phase 2 循环框
\node[loop_box, minimum width=6.5cm, minimum height=4.2cm] (p2_loop_box) at (4.5, -7.8) {};
\node[anchor=north, font=\footnotesize\bfseries] at (4.5, -5.8) {PHASE2\_LOOP};

\node[state_phase2] (p2_check) at (4.5, -6.8) {SOB 判定\\(state, cam)有效?};
\node[state_phase2] (coord_trans) at (2.5, -8.2) {COORD\\TRANSFORM};
\node[state_phase2] (residual) at (4.5, -8.2) {RESIDUAL\\CALC};
\node[state_phase2] (jacobian) at (6.5, -8.2) {JACOBIAN\\CALC};
\node[state_phase2] (hessian) at (4.5, -9.4) {HESSIAN\\ACCUM};

\node[state_phase3] (p3_init) at (0, -11) {PHASE3\_INIT\\初始化Schur循环};

% Phase 3 循环框
\node[loop_box, minimum width=6.5cm, minimum height=3.5cm] (p3_loop_box) at (4.5, -12.8) {};
\node[anchor=north, font=\footnotesize\bfseries] at (4.5, -10.9) {PHASE3\_LOOP};

\node[state_phase3] (v_inv) at (2.5, -12) {$\mathbf{H}_{mm}^{-1}$计算\\(j==0)};
\node[state_phase3] (wvinv) at (6.5, -12) {$\mathbf{H}_{pm}\mathbf{H}_{mm}^{-1}$\\计算};
\node[state_phase3] (schur_update) at (4.5, -13.5) {Schur 更新\\$\mathbf{H}_{pp}, \mathbf{b}_p$};

\node[state_write] (write_out) at (0, -15) {WRITE\_OUTPUT\\DMA写回结果};

\node[state_done] (complete) at (0, -16.8) {COMPLETE\\置位ICR.COMPLETE};

% ========== 状态转移 ==========
% IDLE -> LOAD_PARAMS
\draw[arrow] (idle) -- node[right, condition] {IP\_ENABLE\\$\land$ START\_H} (load_param);

% LOAD_PARAMS -> LOAD_INPUTS
\draw[arrow] (load_param) -- node[right, condition] {读取完成} (load_input);

% LOAD_INPUTS -> PHASE2_INIT
\draw[arrow] (load_input) -- node[right, condition] {加载完成} (p2_init);

% PHASE2_INIT -> PHASE2_LOOP
\draw[arrow] (p2_init) -- ++(2, 0) |- (p2_check);

% Phase 2 内部流程
\draw[arrow] (p2_check) -- node[left, condition] {有效} (coord_trans);
\draw[arrow] (coord_trans) -- (residual);
\draw[arrow] (residual) -- (jacobian);
\draw[arrow] (jacobian) -- (hessian);

% Phase 2 跳过路径
\draw[arrow, dashed] (p2_check.east) -- ++(1.5, 0) node[right, condition] {无效\\跳过} |- (hessian.east);

% Phase 2 循环返回
\draw[arrow] (hessian.south) -- ++(0, -0.3) -- ++(3.5, 0) -- ++(0, 3.0) -- ++(-3.5, 0) -- (p2_check.east);
\node[condition] at (8.5, -8) {cam++\\state++\\pt++};

% PHASE2 -> PHASE3_INIT
\draw[arrow] (hessian.west) -- ++(-3, 0) |- node[left, condition, pos=0.75] {pt $\geq$ num\_pts} (p3_init);

% PHASE3_INIT -> PHASE3_LOOP
\draw[arrow] (p3_init) -- ++(2, 0) |- (v_inv);

% Phase 3 内部流程
\draw[arrow] (v_inv) -- (wvinv);
\draw[arrow] (wvinv) -- (schur_update);

% Phase 3 循环返回
\draw[arrow] (schur_update.south) -- ++(0, -0.3) -- ++(3.5, 0) -- ++(0, 2.0) -- ++(-5.5, 0) -- (v_inv.west);
\node[condition] at (8.5, -12.5) {pt++\\j++};

% PHASE3 -> WRITE_OUTPUT
\draw[arrow] (schur_update.west) -- ++(-3, 0) |- node[left, condition, pos=0.8] {j $\geq$ state\_size} (write_out);

% WRITE_OUTPUT -> COMPLETE
\draw[arrow] (write_out) -- node[right, condition] {写入完成} (complete);

% COMPLETE -> IDLE (返回)
\draw[arrow] (complete.west) -- ++(-2, 0) -- ++(0, 16.8) -- (idle.west);

\end{tikzpicture}

% 字体: 中文使用宋体（继承文档设置），英文使用 Times New Roman（newtxtext）

\begin{tikzpicture}[
    every node/.append style={font=\rmfamily},  % 确保使用衬线字体（Times/宋体）
    state/.style={draw, thick, rounded corners=5pt, minimum width=2.6cm, minimum height=0.9cm, align=center, font=\footnotesize},
    state_idle/.style={state, fill=gray!20},
    state_load/.style={state, fill=blue!15},
    state_phase2/.style={state, fill=green!15},
    state_phase3/.style={state, fill=orange!15},
    state_write/.style={state, fill=purple!15},
    state_done/.style={state, fill=yellow!20},
    arrow/.style={->, >=Stealth, thick},
    condition/.style={font=\tiny, align=center},
    loop_box/.style={draw, dashed, thick, rounded corners=8pt, fill=white, fill opacity=0.7},
]

% ========== 状态节点 ==========
\node[state_idle] (idle) at (0, 0) {IDLE};

\node[state_load] (load_param) at (0, -1.8) {LOAD\_PARAMS\\读取参数};

\node[state_load] (load_input) at (0, -3.6) {LOAD\_INPUTS\\读取SOB/观测/路标};

\node[state_phase2] (p2_init) at (0, -5.4) {PHASE2\_INIT\\初始化循环计数器};

% Phase 2 循环框
\node[loop_box, minimum width=6.5cm, minimum height=4.2cm] (p2_loop_box) at (4.5, -7.8) {};
\node[anchor=north, font=\footnotesize\bfseries] at (4.5, -5.8) {PHASE2\_LOOP};

\node[state_phase2] (p2_check) at (4.5, -6.8) {SOB 判定\\(state, cam)有效?};
\node[state_phase2] (coord_trans) at (2.5, -8.2) {COORD\\TRANSFORM};
\node[state_phase2] (residual) at (4.5, -8.2) {RESIDUAL\\CALC};
\node[state_phase2] (jacobian) at (6.5, -8.2) {JACOBIAN\\CALC};
\node[state_phase2] (hessian) at (4.5, -9.4) {HESSIAN\\ACCUM};

\node[state_phase3] (p3_init) at (0, -11) {PHASE3\_INIT\\初始化Schur循环};

% Phase 3 循环框
\node[loop_box, minimum width=6.5cm, minimum height=3.5cm] (p3_loop_box) at (4.5, -12.8) {};
\node[anchor=north, font=\footnotesize\bfseries] at (4.5, -10.9) {PHASE3\_LOOP};

\node[state_phase3] (v_inv) at (2.5, -12) {$\mathbf{H}_{mm}^{-1}$计算\\(j==0)};
\node[state_phase3] (wvinv) at (6.5, -12) {$\mathbf{H}_{pm}\mathbf{H}_{mm}^{-1}$\\计算};
\node[state_phase3] (schur_update) at (4.5, -13.5) {Schur 更新\\$\mathbf{H}_{pp}, \mathbf{b}_p$};

\node[state_write] (write_out) at (0, -15) {WRITE\_OUTPUT\\DMA写回结果};

\node[state_done] (complete) at (0, -16.8) {COMPLETE\\置位ICR.COMPLETE};

% ========== 状态转移 ==========
% IDLE -> LOAD_PARAMS
\draw[arrow] (idle) -- node[right, condition] {IP\_ENABLE\\$\land$ START\_H} (load_param);

% LOAD_PARAMS -> LOAD_INPUTS
\draw[arrow] (load_param) -- node[right, condition] {读取完成} (load_input);

% LOAD_INPUTS -> PHASE2_INIT
\draw[arrow] (load_input) -- node[right, condition] {加载完成} (p2_init);

% PHASE2_INIT -> PHASE2_LOOP
\draw[arrow] (p2_init) -- ++(2, 0) |- (p2_check);

% Phase 2 内部流程
\draw[arrow] (p2_check) -- node[left, condition] {有效} (coord_trans);
\draw[arrow] (coord_trans) -- (residual);
\draw[arrow] (residual) -- (jacobian);
\draw[arrow] (jacobian) -- (hessian);

% Phase 2 跳过路径
\draw[arrow, dashed] (p2_check.east) -- ++(1.5, 0) node[right, condition] {无效\\跳过} |- (hessian.east);

% Phase 2 循环返回
\draw[arrow] (hessian.south) -- ++(0, -0.3) -- ++(3.5, 0) -- ++(0, 3.0) -- ++(-3.5, 0) -- (p2_check.east);
\node[condition] at (8.5, -8) {cam++\\state++\\pt++};

% PHASE2 -> PHASE3_INIT
\draw[arrow] (hessian.west) -- ++(-3, 0) |- node[left, condition, pos=0.75] {pt $\geq$ num\_pts} (p3_init);

% PHASE3_INIT -> PHASE3_LOOP
\draw[arrow] (p3_init) -- ++(2, 0) |- (v_inv);

% Phase 3 内部流程
\draw[arrow] (v_inv) -- (wvinv);
\draw[arrow] (wvinv) -- (schur_update);

% Phase 3 循环返回
\draw[arrow] (schur_update.south) -- ++(0, -0.3) -- ++(3.5, 0) -- ++(0, 2.0) -- ++(-5.5, 0) -- (v_inv.west);
\node[condition] at (8.5, -12.5) {pt++\\j++};

% PHASE3 -> WRITE_OUTPUT
\draw[arrow] (schur_update.west) -- ++(-3, 0) |- node[left, condition, pos=0.8] {j $\geq$ state\_size} (write_out);

% WRITE_OUTPUT -> COMPLETE
\draw[arrow] (write_out) -- node[right, condition] {写入完成} (complete);

% COMPLETE -> IDLE (返回)
\draw[arrow] (complete.west) -- ++(-2, 0) -- ++(0, 16.8) -- (idle.west);

\end{tikzpicture}
